\section{Conclusion}
\label{sec:conclusion}

We have presented a study of antisemitism detection on the GoldStandard2024 dataset along three connected axes --- accuracy on the positive class, calibration per sub-group, and counterfactual stability under identity-token swaps --- and proposed Counterfactual Calibration Invariance (CCI v2), a training-time loss penalising the Jensen-Shannon divergence between predictions on a sentence and on its identity-token-swapped counterfactual within a religiously-symmetric vocabulary. A short structural lemma rules out scalar post-hoc Temperature Scaling as a means of reducing the hard counterfactual flip rate at the decision threshold $1/2$; we verify the lemma bit-exactly on six Model~C seeds as a regression test.

Empirically, CCI v2 reduces the hard counterfactual flip rate by 49\% on DeBERTa-v3-base (184M), 55\% on DeBERTa-v3-large (440M, same family), and 73\% on RoBERTa-large (355M, different family) versus a strong focal+identity-masking baseline at iso-architecture, with positive-class F\textsubscript{1} preserved or improved at scale (0.72 best). A two-family paired-bootstrap suite with Holm-Bonferroni correction and exact sign-test backup (54 tests in the main family, 6 in the within-architecture transfer family) supports the CFR claims under multiple-comparison correction. Two further findings round out the empirical picture: (i)~the headline observation: post-hoc multicalibration, while improving global subgroup ECE, \emph{increases} the L\textsubscript{1}-norm soft counterfactual flip rate by 96\% on this task --- a tension between two well-respected fairness desiderata that, to our reading, has not been quantitatively documented for text classifiers; (ii)~the FNED trade-off is mild: CCI v2 (fixed T, JS) shows only $+2.9\%$ versus the strongest CLP baseline (Davani $\lambda{=}0.5$) at base scale, and \emph{improves} FNED versus same-architecture Model~C ($-7.0\%$ on base, $-1.5\%$ on DeBERTa-v3-large). The substantial regression we initially documented turned out to be specific to a different (learned-$T$) CCI v2 cell rather than to the fixed-$T$ pivot we now report. We additionally document a bimodal ECE distribution under focal loss with F\textsubscript{1}-based early stopping as a supporting observation, traced mechanistically to the F\textsubscript{1}-peak epoch.

\paragraph{Practical implications.}
For content moderation at scale, training-time counterfactual invariance is a low-cost addition to existing pipelines: the swap engine is purely algorithmic, the loss adds one paired forward pass per training example, and the resulting model has no inference-time overhead. The structural Lemma~\ref{thm:ts-cfr} cautions against the assumption that post-hoc fairness fixes (Temperature Scaling, multicalibration) are universal: they optimise specific objectives and may degrade others. Practitioners with both subgroup-conditional calibration and counterfactual-invariance requirements should adopt a training-time intervention rather than stack post-hoc tools.

\paragraph{Future work.}
\begin{enumerate}[leftmargin=*,itemsep=2pt]
    \item \textbf{A tight Lipschitz-style certificate for CFR\textsubscript{hard}.} Lemma~\ref{thm:ts-cfr} characterises what TS cannot do; a corresponding upper bound of the form
    $\mathrm{CFR}_{\text{hard}} \le \mathbb{E}_x[L_{\text{local}}(f; v_g) \cdot \|\Delta e\|_2 / m(x)]$
    in terms of the classifier's local Lipschitz constant along the identity-token direction would convert our empirical reduction into a guarantee.
    \item \textbf{LLM-likelihood-based symmetry filter.} Our heuristic symmetry classifier is rule-based; a language-model-likelihood scorer (already plumbed into the pipeline) could improve the swap manifest's coverage on borderline cases.
    \item \textbf{Cross-language transfer.} Antisemitism manifests differently across languages, with distinct historical tropes; cross-lingual transfer to Hebrew, Arabic, and major European languages would test the generality of the swap-engine + symmetry-filter formulation.
    \item \textbf{Cross-identity transfer.} The method is identity-agnostic in principle (the symmetric vocabulary is the only identity-specific component); empirical validation on anti-Muslim, anti-Black, or anti-LGBTQ+ hate-speech tasks would establish the method's generality.
    \item \textbf{Multi-label severity-graded classification.} Moving beyond binary detection to classify the \emph{type} of antisemitic expression (conspiracy, dehumanization, Holocaust distortion, dual loyalty) would provide more actionable information for moderators and richer signal for the counterfactual-invariance objective.
\end{enumerate}

Automated antisemitism detection is not a solved problem, but the three-axis evaluation framework, the multicalibration-vs-counterfactual-invariance tension we surface, and the architecture-and-family transfer of CCI v2 all suggest that progress requires moving beyond accuracy as the sole optimisation target. We hope the publicly released training code, multi-seed checkpoints, and Bonferroni-corrected significance suite will be useful to the community working on this and adjacent problems.
